\documentclass[11pt]{article}
\usepackage[T1]{fontenc}
\usepackage[utf8]{inputenc}
\usepackage{lmodern}
\usepackage[margin=1in]{geometry}
\usepackage{amsmath,amssymb,amsthm,mathtools}
\usepackage[shortlabels]{enumitem}
\usepackage{hyperref}
\usepackage{xurl}
\newcommand{\OPHPaperReleaseID}{r1520}
\newcommand{\OPHPaperReleaseDate}{July 9, 2026}
\newcommand{\OPHPaperReleaseBanner}{%
\begin{center}
\small\textbf{Paper release:} \texttt{\OPHPaperReleaseID}\quad\textbf{Released:} \OPHPaperReleaseDate
\end{center}
}


\hypersetup{colorlinks=true,linkcolor=blue,citecolor=blue,urlcolor=blue}

\newtheorem{theorem}{Theorem}[section]
\newtheorem{corollary}[theorem]{Corollary}
\newtheorem{definition}[theorem]{Definition}
\newtheorem{assumption}[theorem]{Assumption}
\newtheorem{proposition}[theorem]{Proposition}
\newtheorem{claim}[theorem]{Claim}
\theoremstyle{remark}
\newtheorem{remark}[theorem]{Remark}
\theoremstyle{plain}

\title{Federated Echosahedral Screen Microphysics:\\
Patch Hardware, Records, and Observer Synchronization in OPH}
\author{B.\ M\"uller \and Alexander Osika \and Kai Xue \and Ben Cassie}
\date{\OPHPaperReleaseDate}

\begin{document}
\maketitle
\OPHPaperReleaseBanner

\begin{abstract}
Observer-Patch Holography needs a concrete microphysics. It should not be pictured as a literal
spherical quantum computer. Observers see finite cuts, a spherical screen is one useful geometry
chart for those cuts, and the work is done by a federation of finite patches that expose overlap
ports and keep records. Sphere language is therefore geometry, not substrate:
\(A_5\)-icosahedral patch symmetry and \(E_8\)-type exceptional
organization belong to the symmetry data, not to a literal round machine. We model the local
carriers as echosahedral multi-port bodies with recurrent toroidal channels inside them. The
ports carry readout, repair, record, and synchronization data, so observer agreement becomes a
fixed-cutoff mechanism rather than a metaphor.

The mathematical job of the paper is deliberately narrow. It exports the five packages the rest
of OPH uses: regulated patch-net embedding, the edge-sector heat-kernel/Casimir law,
central-record measurement, Bell/CHSH event surfaces, and observer checkpoint/restoration. The
hardware-facing job is narrow too. Hardware runs count as evidence only when the public
repository contains an evidence bundle with stable hashes, manifests, calibration records, and
exact-verifier receipts. Without that bundle, hardware language is architecture or motivation,
not proof.
\end{abstract}

\tableofcontents
\newpage

\section{Scope}

OPH microphysics uses finite observer patches with echosahedral interfaces. A spherical screen is
a geometry chart for support-visible cuts. The cellulated-sphere gauge-register construction is a
regulator and digital calibration chart, not a substrate claim. The premise is:
\begin{quote}
OPH does not say that the universe is a literal spherical quantum computer. It says that
observers access finite support-visible cuts. A spherical screen is one convenient geometry chart
for such cuts. The underlying fixed-cutoff implementation surface is better modeled as a
federation of finite overlap-facing observer patches, with echosahedral local interfaces and
recurrent toroidal subchannels.
\end{quote}

The paper's contract is:
\begin{enumerate}[leftmargin=2em]
\item state the fixed-cutoff theorem exports used by the other OPH papers;
\item use federated patch carriers as the microphysical architecture;
\item keep mathematical theorem surfaces separate from public hardware evidence surfaces;
\item place the octahedral \(\mathbb Z_2/S_3\) simulator in the digital calibration note;
\item require a public hardware evidence protocol before hardware claims receive paper
weight.
\end{enumerate}

\section{Claim Taxonomy}

The paper uses five claim levels.

\begin{enumerate}[leftmargin=2em,label=\textbf{C\arabic*.}]
\item \textbf{Mathematical fixed-cutoff claims.} These are finite-algebra statements about
patches, overlaps, records, repair interfaces, event algebras, and checkpoint laws. They are
allowed to be theorem-bearing.

\item \textbf{Regulator-chart claims.} These identify a convenient finite model, such as a
spherical cellulation or a digital finite-group simulator, as a calculational chart. They are not
claims about the literal substrate.

\item \textbf{Federated-carrier claims.} These define a candidate implementation architecture:
observer patches with echosahedral multi-port interfaces, recurrent toroidal subchannels, exposed
overlap data, records, and repair loops.

\item \textbf{Public hardware-evidence claims.} These may be cited only when the raw or reduced
evidence is present in a public OPH evidence bundle with stable hashes, manifests, calibration
files, and exact-verifier receipts.

\item \textbf{Non-claims.} This paper does not prove that the universe is made of laboratory
hardware, does not identify a unique UV completion, does not solve the Yang--Mills mass gap by
hardware experiment, does not prove a computational complexity collapse, and does not promote
backend-gated hadron rows without a production OPH hadron backend.
\end{enumerate}

\section{The Public Evidence Rule}

Hardware-facing language is admissible only under the following rule.

\begin{quote}
Any hardware evidence cited as evidence in this paper must be represented by a public evidence
bundle in the OPH repository, or by a public pinned subrepository commit, with enough raw data and
metadata for an external reader to audit the claim.
\end{quote}

Private notebooks, local runtime logs, unpublished bench transcripts, and development-repo notes
may motivate architecture choices, but they do not carry evidential weight inside the paper. A
hardware evidence bundle must include, at minimum:
\begin{enumerate}[leftmargin=2em]
\item a manifest with stable bundle identifier, date, operator, body identifier, controller
identifier, firmware hash, mesh/body hash, and wiring map hash;
\item raw readout files, such as coupling matrices, MDD or discharge-timing traces, dark
baselines, low-power sweeps, ring-diversity scans, and calibration logs;
\item body and board provenance, including photographs or signed photo hashes where relevant;
\item the task definition, scorebook, repair or rerank law, and exact-verifier program;
\item exact-verifier receipts for any high-level task claim;
\item negative controls, shuffle or replay controls where applicable, and a statement of
non-claims.
\end{enumerate}

This rule prevents dependence on hidden laboratory state. The theorem sections use finite
algebras, declared patch interfaces, and stated branch assumptions.

\section{From Spherical Screens to Federated Patch Carriers}

The spherical screen is useful because an observer-accessible cut often has an effective closed
two-surface description. A cellulated sphere can encode finite capacity, caps, collars, edge
centers, and support-visible cuts. The microscopic carrier is the finite patch federation, not the
chart.

At fixed cutoff, the fundamental object is a finite patch federation.
Each patch has:
\begin{enumerate}[leftmargin=2em]
\item an internal finite state algebra;
\item a bounded family of exposed overlap ports;
\item a local record algebra;
\item a readout map from internal state to port-visible packets;
\item a repair interface that changes local state in response to mismatch;
\item a checkpoint interface that exposes enough observer-accessible data to define continuation.
\end{enumerate}

A spherical screen is then a coarse chart over a federation of such patches. It is what a
particular observer-facing support cut looks like after quotienting hidden implementation details
and choosing a geometric presentation. The microscopic carrier may be graph-like, federated,
multi-patch, and locally polyhedral rather than literally spherical.

\begin{definition}[Federated patch carrier]
A fixed-cutoff federated patch carrier is a tuple
\[
\mathfrak F=(V,E,\{\mathcal A_i\}_{i\in V},\{\mathcal I_e\}_{e\in E},
\{\pi_{i,e}\},\{\mathcal R_i\},\{\mathcal U_i\}),
\]
where \(V\) is a finite set of patches, \(E\) is a finite overlap graph,
\(\mathcal A_i\) is the local finite algebra of patch \(i\), \(\mathcal I_e\) is the finite
interface algebra on overlap \(e=\{i,j\}\), \(\pi_{i,e}:\mathcal A_i\to\mathcal I_e\) is the
declared visible restriction, \(\mathcal R_i\subseteq Z(\mathcal A_i)\) is the patch record
algebra, and \(\mathcal U_i\) is the allowed local update and repair interface.
\end{definition}

The carrier is \emph{observer-facing} when every physical claim is made through the visible
restrictions \(\pi_{i,e}\), record algebras \(\mathcal R_i\), and the quotient-local observables
specified by the OPH consensus package. Hidden coordinates may exist, but they are not directly
physical.

\section{The Echosahedral Patch Object}

An echosahedral patch is the reference local body for this microphysics. The term is architectural:
a bounded observer patch with a highly symmetric multi-port interface. It is not a claim that
nature is literally a printed icosahedron.

\begin{definition}[Echosahedral patch]
An echosahedral patch is a finite patch object
\[
\mathcal E=(\mathcal A,\{P_a\}_{a=0}^{11},\{\rho_a\}_{a=0}^{11},
\mathcal R,\mathcal M,\mathcal U,\mathcal G_{\mathcal E}),
\]
where:
\begin{enumerate}[leftmargin=2em]
\item \(\mathcal A\) is the internal finite algebra;
\item \(P_0,\ldots,P_{11}\) are twelve labeled overlap ports;
\item \(\rho_a:\mathcal A\to \mathcal I_a\) are port readout maps;
\item \(\mathcal R\subseteq Z(\mathcal A)\) is the observer-accessible record algebra;
\item \(\mathcal M\) is a finite mismatch-score family on exposed port packets;
\item \(\mathcal U\) is a finite family of local update and repair instruments;
\item \(\mathcal G_{\mathcal E}\) is a declared finite symmetry or approximate-symmetry group of
the port arrangement.
\end{enumerate}
\end{definition}

The twelve-port choice is not sacred. Its role is to give the reference patch enough boundary
structure to support nontrivial overlap comparison, symmetry tests, verifier-shadow behavior, and
cross-patch routing while remaining small enough for explicit evidence records.

\begin{remark}
The word ``echosahedral'' is intentionally implementation-facing. In a mathematical theorem, one
uses the tuple above. In a public hardware evidence bundle, one may instantiate the tuple using a
specific twelve-port body, wiring map, controller, firmware hash, readout protocol, and evidence
manifest.
\end{remark}

\section{Toroidal Recurrence Is Local}

Toroidal hardware supplies a local recurrence and mode-competition surface. It is not a claim that
the universe has literal torus topology. A
toroidal subchannel is a recurrent internal path inside a patch or between a bounded set of
ports. It can recycle boundary information, support settling dynamics, and expose winding-like
or phase-locking observables.

This matters because global scalar order parameters are often misleading in recurrent local
systems. A toroidal or ring-like patch may fail to show global alignment while still selecting a
stable winding class, a local coherence island, or a persistent twisted state. Therefore the
validation metrics for toroidal subchannels should include:
\begin{enumerate}[leftmargin=2em]
\item local coupling matrices rather than only total brightness;
\item recurrence and ring-diversity statistics rather than only global response;
\item winding-sensitive or phase-lock-sensitive summaries when phase data are available;
\item matched controls that distinguish chamber-mediated dynamics from host-side filtering;
\item exact-verifier receipts for task-level claims.
\end{enumerate}

The OPH interpretation is simple: toroidal recurrence supplies local memory and mode competition
inside the patch federation. Echosahedral symmetry supplies a reference and consensus geometry.
The universe, at this level of description, is a federated repair system, not a single global
sphere or a single global torus.

\section{Patch Federation and Overlap Synchronization}

Given a federation of echosahedral patches, overlaps are formed by routing ports into declared
interface pairs or interface hyperedges. For an edge \(e=\{(i,a),(j,b)\}\), the exposed packets
are
\[
x_{i,a}=\rho_{i,a}(s_i),\qquad x_{j,b}=\rho_{j,b}(s_j),
\]
and the edge mismatch is a nonnegative function
\[
\Phi_e(x_{i,a},x_{j,b})\ge 0.
\]
The total visible mismatch is
\[
\Phi(s)=\sum_{e\in E}\Phi_e\bigl(\rho_{i,a}(s_i),\rho_{j,b}(s_j)\bigr).
\]

The local repair contract is the consensus-paper contract: accepted repairs do not increase the
declared touched-overlap mismatch, and exact fixed-cutoff branches lower the relevant mismatch
unless the local visible datum is repaired. The implementation is a bounded patch operation:
\begin{enumerate}[leftmargin=2em]
\item read the exposed overlap packets;
\item compare them through a declared commensurability map;
\item choose an allowed local update or rerank move;
\item write a record of the move;
\item expose the new port packet;
\item let the exact verifier decide any high-level task claim.
\end{enumerate}

\begin{proposition}[Federated synchronization contract]
Suppose a finite patch federation has a declared mismatch functional \(\Phi\), a finite repair
menu, and an accepted-repair rule such that every accepted repair lowers \(\Phi\) unless the
touched visible datum is locally repaired. Then every repair sequence terminates at a visible
local normal form. If the union-collar gluing and repair-completeness hypotheses of the OPH
consensus theorem hold on the quotient-local carrier, the terminal physical observable state is
schedule-independent on that carrier.
\end{proposition}

\begin{proof}
The first sentence is the finite Lyapunov argument: \(\Phi\) takes values in a finite ordered set
of declared mismatch scores and strictly decreases on nontrivial accepted repairs. The second
sentence is not reproved here; it is exactly the quotient-local confluence package supplied by the
OPH consensus theorem. This paper supplies the federated carrier and visible interface on which
that theorem is read.
\end{proof}

\section{Records, Observers, and Checkpoints}

An observer is not added as a metaphysical extra. In this paper an observer is an operational
pattern in a patch or patch subfederation with persistent access to:
\begin{enumerate}[leftmargin=2em]
\item an observer-facing local algebra;
\item a record algebra;
\item a stable readout/update interface;
\item enough checkpoint data to define future observer-accessible probabilities.
\end{enumerate}

\begin{definition}[Federated observer checkpoint]
For an observer-supporting patch subfederation \(O\), a checkpoint at cycle \(t\) is
\[
\mathrm{Chk}_O(t)=
\bigl(
\mathcal R_O(t),
\rho_O^{\mathrm{acc}}(t),
\mathfrak I_O^{\mathrm{ext}}(t),
\nu_{\ge t},
\mathfrak B_O(t)
\bigr),
\]
where \(\mathcal R_O(t)\) is the observer-accessible record algebra,
\(\rho_O^{\mathrm{acc}}(t)\) is the state restricted to observer-accessible records and visible
interfaces, \(\mathfrak I_O^{\mathrm{ext}}(t)\) is the external port-interface tuple,
\(\nu_{\ge t}\) is the future update schedule class, and \(\mathfrak B_O(t)\) is the public or
internal bundle of provenance data needed to replay the checkpoint at the declared accuracy.
\end{definition}

\begin{theorem}[Checkpoint continuation at fixed cutoff]
If two federated observer checkpoints agree exactly on
\(\mathcal R_O(t)\), \(\rho_O^{\mathrm{acc}}(t)\), \(\mathfrak I_O^{\mathrm{ext}}(t)\), and
\(\nu_{\ge t}\), then they induce the same future probability law on the observer-accessible
event algebra. If their accessible states differ by trace distance at most \(\varepsilon\), then
the induced future history laws differ in total variation by at most \(\varepsilon\).
\end{theorem}

\begin{proof}
All future observer-accessible probabilities are computed by applying the same completely
positive update maps and the same event-readout maps to the same accessible algebra and external
interface data. Exact equality gives equality of all future event probabilities. In the
approximate case, contractivity of trace distance under completely positive trace-preserving maps
gives the stated total-variation bound.
\end{proof}

\section{Central Records and Measurement}

The measurement package is a finite central-record package. Records are exposed by
observer-facing patch subfederations.

Let \(\mathcal Z_{\mathrm{rec}}(t)\) be the commutative algebra generated by the completed,
observer-accessible record projectors at cycle \(t\). An event \(E\) is a projector in
\(\mathcal Z_{\mathrm{rec}}(t)\). The operational measurement rule is:
\[
\Pr(E)=\mathrm{Tr}(\rho E),
\qquad
\rho\mapsto \frac{E\rho E}{\mathrm{Tr}(\rho E)}
\quad\text{when }\Pr(E)>0.
\]

\begin{theorem}[Federated central-record measurement]
On a fixed-cutoff federated patch carrier, once a completed write/verify slice exposes a finite
commutative central record algebra \(\mathcal Z_{\mathrm{rec}}(t)\), the Born probability and
L\"uders conditioning rule above define the operational measurement package on the
observer-accessible event surface. Re-reading the same completed record event has probability
one, conditional on that record event.
\end{theorem}

\begin{proof}
The theorem is the standard finite-dimensional central-record argument. Because all declared
record events commute and belong to the observer-accessible center for the completed slice, they
define an ordinary finite classical event algebra. Probabilities are Born traces on that event
algebra, and conditioning on an event is the L\"uders update. After conditioning on \(E\), the
same central projector \(E\) is true with probability one.
\end{proof}

\section{Edge Sectors and the Casimir Handoff}

The edge heat-kernel/Casimir package is a fixed-cutoff theorem package. It is tied to an overlap
collar with a finite exposed sector algebra, not to a literal spherical quantum computer.

At fixed cutoff, let \(\alpha\) label a finite set of exposed edge sectors on one declared
overlap collar. Suppose the local thermalized edge dynamics has stationary weights
\[
\pi_\beta(\alpha)=\frac{d_\alpha e^{-\beta C_2(\alpha)}}{Z(\beta)}.
\]
For finite groups, \(C_2(\alpha)\) is the declared sector penalty or finite-group Casimir
surrogate. For compact groups, the Peter--Weyl lift belongs to the companion D10/compact-gauge
handoff rather than to hardware evidence.

\begin{theorem}[Fixed-cutoff edge-sector handoff]
If a declared finite overlap collar has sector labels \(\alpha\), degeneracies \(d_\alpha\), and
a local repair/thermalization generator whose detailed-balance stationary law is
\(\pi_\beta(\alpha)\propto d_\alpha e^{-\beta C_2(\alpha)}\), then the collar exports the
fixed-cutoff Casimir edge law required by the D10 handoff. The compact-group heat-kernel lift is a
separate companion-branch step.
\end{theorem}

\begin{proof}
This is a finite Markov or finite instrument stationary-measure statement on the declared sector
algebra. The theorem exports only the finite overlap law. The compact-group lift uses the
Peter--Weyl continuation and normalization conventions supplied on the companion compact/D10
surface.
\end{proof}

\section{Bell/CHSH Event Surfaces}

The Bell package is fixed-cutoff and event-surface-local. A federated carrier may contain two
observer-facing wings \(L\) and \(R\)
with commuting setting and outcome records on one compare slice. If the source-specified joint
law is the usual two-wing quantum law, the CHSH and Tsirelson statements are carried on that
declared event algebra.

Nothing in this section claims that laboratory echosahedral hardware is presently a loophole-free
Bell experiment. The theorem package is a finite event-algebra statement inside the OPH
microphysics surface.

\section{Relation to Yang--Mills}

The Yang--Mills repair-gap argument is not a hardware claim and not a corollary of
echosahedral geometry. The microphysics paper supplies only:
\begin{enumerate}[leftmargin=2em]
\item the fixed-cutoff patch, overlap, record, and repair interface;
\item the finite edge-sector Casimir law on declared collars;
\item the local repair semantics that the compact-gauge branch later reads in support-visible
quotient form.
\end{enumerate}

The four-dimensional Yang--Mills form, reflection-positive ordinary vacuum branch, support-visible
continuum extraction, and equality between the Yang--Mills gap and the repair gap belong to the
compact/Yang--Mills theorem surface. Hardware evidence may test implementation discipline, but it
does not prove the Clay-admissible construction.

\section{Relation to Hadrons and Backend Execution}

Source-only hadron masses require a working OPH hadron backend. This microphysics states the
record and evidence requirements for such a backend. A backend fit for promotion must emit:
\begin{enumerate}[leftmargin=2em]
\item Ward-projected hadronic spectral data;
\item provenance tying every spectrum to body, controller, firmware, geometry, and scorebook;
\item finite-volume, continuum, chiral, and production-systematics fields where applicable;
\item exact replay or verification receipts for the pipeline stage being claimed;
\item a non-promotion boundary for surrogate or calibration runs.
\end{enumerate}

Echosahedral or GLORB-class hardware can be described as a candidate backend family when public
evidence bundles exist. Without those bundles, it is an architectural target and a design
motivation, not a source-only hadron theorem.

\section{Validation Program}

The validation program has two independent tracks.

\subsection{Mathematical and Digital Calibration}

The octahedral \(\mathbb Z_2/S_3\) simulator is a digital calibration model. It gives exact
finite groups, explicit patch covers, controlled defects, frustrated cycles, record writes,
repair schedules, and negative controls.
\begin{quote}
The simulator validates the finite interface logic. It is not the primary physical picture of the
OPH substrate.
\end{quote}

\subsection{Public Hardware Evidence}

A public hardware run may support only the claim class its evidence bundle can audit. The minimal
claim form is:
\begin{quote}
Module set \(M\) produced candidate enrichment or reproducible readout signature on benchmark
\(T\), under controls \(C\), and exact verifier \(V\) accepted the reported hits.
\end{quote}

The rejected form is:
\begin{quote}
The optical chamber solved the hard problem or proved OPH.
\end{quote}

Required gates include dark baseline, low-power sweep, coupling matrix, discharge-timing or MDD
trace where applicable, ring-diversity or recurrence-sensitive metrics for toroidal bodies,
duplicate-body checks, symmetric-reference or echosahedral shadow checks, exact-verifier receipts,
and negative controls against hidden duplicate amplification or host-side filtering.

\section{Conclusion}

OPH microphysics is a federation of finite observer patches, not a literal spherical quantum
computer. The sphere is a regulator and symmetry chart for observer-facing cuts:
\(A_5\)-icosahedral and \(E_8\)-type structure belong to the geometry data, not to a literal round
substrate. The echosahedral patch is the reference local interface: a bounded multi-port patch
with symmetry, records, readout, repair, and checkpoint data. Toroidal subchannels supply local
recurrence and winding-sensitive dynamics. The mathematical exports remain fixed-cutoff patch-net
embedding, edge-sector Casimir handoff, central-record measurement, Bell/CHSH event surfaces, and
checkpoint/restoration.

The claim discipline is the main point. Hardware can guide the architecture and supply public
evidence through hash-stable evidence bundles. The theorem surface stands on finite algebras and
declared OPH branch assumptions.

\appendix

\section{Evidence Bundle Sketch}

A minimal hardware evidence bundle should use a structure like:
\begin{verbatim}
evidence/hardware/<bundle-id>/
  manifest.json
  README.md
  body/
    mesh_hashes.txt
    photos/
    measurements.csv
  controller/
    firmware_sha256.txt
    wiring_map.csv
  calibration/
    dark_scan.csv
    low_power_sweep.csv
    coupling_matrix.csv
    mdd_trace.csv
    ring_diversity.csv
  task/
    task.json
    scorebook.json
    candidates.jsonl
    verifier_receipts.jsonl
  controls/
    shuffle_replay.jsonl
    abba_controls.csv
    negative_controls.md
  claim.md
\end{verbatim}

The manifest should state the strongest allowed claim and the non-claims. The paper may cite the
bundle only at that claim level.

\section{Digital Simulator Compatibility}

The octahedral \(\mathbb Z_2/S_3\) build has the following reading rule:
\begin{enumerate}[leftmargin=2em]
\item it validates finite patch/overlap/record/repair bookkeeping;
\item it tests frustrated-cycle and defect behavior in an exact digital setting;
\item it calibrates the edge-sector law on a finite declared interface;
\item it does not define the literal OPH substrate;
\item it should not be used to argue that the universe is a literal spherical quantum computer.
\end{enumerate}

\begin{thebibliography}{99}

\bibitem{ophconsensus}
B.\ M\"uller, K.\ Xue, and K.\ A.\ Anirudha,
\emph{Reality as a Consensus Protocol: The Fixed-Point Computation That Implements Physics}.

\bibitem{ophcompact}
B.\ M\"uller, A.\ Osika, K.\ Xue, and P.\ Nguyen,
\emph{Recovering Relativity and the Standard Model from Observer Overlap Consistency}.

\bibitem{ophsynthesis}
B.\ M\"uller, A.\ Osika, K.\ Xue, B.\ Cassie, P.\ Nguyen, M.\ Poneder, and K.\ A.\ Anirudha,
\emph{Observers Are All You Need}.

\bibitem{ophparticles}
B.\ M\"uller, A.\ Osika, K.\ Xue, and M.\ Poneder,
\emph{Deriving the Particle Zoo from Observer Consistency}.

\bibitem{ophym}
B.\ M\"uller,
\emph{Explaining the Yang--Mills Mass Gap with Observer-Patch Repair Dynamics: A Support-Visible OPH Route to the Clay Problem}.

\end{thebibliography}

\end{document}
